%% -*- coding: utf-8 -*-
%% EXAMPLE CONTENT — ships with a new SciTeX workspace. Replace it with yours.
\section{Discussion}
\label{sec:discussion}
\pdfbookmark[1]{Discussion}{discussion}

\subsection{What the example shows}

Most of the digit-recognition problem, at this resolution, is linearly
separable in raw pixel space. A rule with no feature engineering behind it
recovers 87.3\% of the held-out labels, and for eight of the ten classes it is
close to the ceiling the sample size allows.

The residual errors are where the example earns its keep. They are not noise
spread evenly across classes; they concentrate on 8 and, to a lesser extent, 9.
A linear decision boundary can only weigh each of the 64 pixels independently.
The digit 8 fills nearly the whole $8\times8$ frame, so it overlaps in lit
pixels with most other digits at once --- which is why its 8 errors scatter
across five different classes (1, 4, 5, 6 and 7) rather than piling onto a
single rival. What separates 8 from those digits is the \emph{arrangement} of
the ink, not which pixels carry it, and arrangement is precisely what a linear
model on raw pixels cannot express. This is the standard motivation for the
kernel and convolutional methods that followed.

\subsection{What it does not show}

Several things, deliberately:

\begin{itemize}
  \item \textbf{This is not a benchmark.} The 500-image subset is a twentieth
    of the source corpus, and the 150-image test set gives each class only 15
    examples. A single reclassified image moves a per-class rate by nearly
    seven points, so the per-digit numbers in Section \ref{sec:results} are
    illustrative, not precise. No confidence intervals are reported because a
    single split cannot support them.
  \item \textbf{Nothing here is new.} The result reproduces textbook behaviour
    on a dataset published in 1998. It is included because it is verifiable,
    not because it is informative.
  \item \textbf{No model comparison was run.} A kernel SVM, a random forest, or
    a small convolutional network would all be expected to do better. None was
    fitted, so no such claim is made.
\end{itemize}

\subsection{Using this project as a starting point}

The structure generalises even though the content does not. Data lives in
\texttt{data/} with its provenance beside it; one seeded script in
\texttt{scripts/} turns that data into everything the manuscript displays; the
manuscript cites the figures rather than restating numbers by hand. Replace the
CSV, adjust the script, and rewrite these sections --- the wiring between them
already works, and the reproduction command stays the same.

%%%% EOF
